\section{The exact force needed to reach a compact based state}
\label{ifr:section}

The annulus defect has a further finite-interval consequence for the
actual signed force.  Retain the entire smooth parent
$(\theta_P,u_P,p_P)$ on $[a,b]$, and let a compact smooth increment
$(h,z,\kappa)$ have $\operatorname{div}z=0$ and $h(a)=z(a)=0$.
The support of the increment is contained in one fixed compact set on
this finite interval.  Write $e^\theta,e^u$ for its complete physical
residual \eqref{ibc:scalar-residual}--\eqref{ibc:vector-residual}.
No assumption that the parent's own force vanishes is made.

For every retained weight $\lambda>0$, define the original-coordinate
quantities
\begin{align}
 E_\lambda(t)^2&=\|z(t)\|_2^2+\lambda\|h(t)\|_2^2,\nonumber\\
 F_\lambda(t)^2&=\|e^u(t)\|_2^2+\lambda\|e^\theta(t)\|_2^2,\nonumber\\
 C_\lambda(t)&=\|\nabla u_P(t)\|_\infty+
             \frac{1+\lambda\|\nabla\theta_P(t)\|_\infty}
                  {2\sqrt\lambda}.
 \label{ifr:weighted-data}
\end{align}
The parent norms may be restricted to the common compact support
neighbourhood; using the full spatial norms also gives the displayed
bounds.  The diffusion coefficient $d>0$ is unchanged.
The exact energy identity is
\begin{align}
 \frac12(E_\lambda^2)'&
   +d\bigl(\|\nabla z\|_2^2+\lambda\|\nabla h\|_2^2\bigr)\nonumber\\
 &= -\int z^T(\nabla u_P)z
    -\lambda\int h\,z\cdot\nabla\theta_P
    +\int h z_2
    +\int z\cdot e^u+\lambda\int h e^\theta .
 \label{ifr:energy-identity}
\end{align}
All integrals use the original Lebesgue measure on $\mathbb R^2$.

\begin{proof}
Multiply the scalar increment equation by $\lambda h$ and the vector
increment equation by $z$, then integrate and add.
Both $u_P$ and $z$ are divergence free.  Therefore
$\int h\,u_P\cdot\nabla h$, $\int h\,z\cdot\nabla h$,
$\int z\cdot(u_P\cdot\nabla)z$, and
$\int z\cdot(z\cdot\nabla)z$ vanish by integration by parts.
The pressure integral $\int z\cdot\nabla\kappa$ also vanishes.
The term $(z\cdot\nabla)u_P$ gives the first right-hand integral.
The scalar parent-gradient term gives the second, and the original
buoyancy $-he_2$ on the residual's left side gives the third with a
positive sign.  Each retained $-d\Delta$ term gives the corresponding
positive gradient norm.  Compact support removes the spatial boundary
terms.  This proves the full identity, including all signs and $\lambda$.
\end{proof}

The parent-gradient and buoyancy terms obey
\[
 \left|-\lambda\int h\,z\cdot\nabla\theta_P+\int h z_2\right|
 \leq\frac{1+\lambda\|\nabla\theta_P\|_\infty}
             {2\sqrt\lambda}E_\lambda^2.
\]
Indeed their sum is bounded by
$(1+\lambda\|\nabla\theta_P\|_\infty)\|h\|_2\|z\|_2$,
and $2\sqrt\lambda\|h\|_2\|z\|_2\leq E_\lambda^2$.
The strain term is at most
$\|\nabla u_P\|_\infty E_\lambda^2$, and the signed forcing
pairing is at most $E_\lambda F_\lambda$.
Dropping only the displayed nonnegative diffusion integral in this
inequality gives the exact finite-interval estimate
\begin{equation}\label{ifr:reachability-bound}
 E_\lambda(t)\leq
 \int_a^t
   \exp\!\left(\int_s^t C_\lambda(\tau)\,d\tau\right)
       F_\lambda(s)\,ds .
\end{equation}
For a proof valid at zeros of $E_\lambda$, put
$E_{\lambda,\varepsilon}=(E_\lambda^2+\varepsilon^2)^{1/2}$.
The energy identity implies
$E_{\lambda,\varepsilon}'\leq
C_\lambda E_{\lambda,\varepsilon}+F_\lambda$.
Multiply by $\exp(-\int_a^t C_\lambda)$, integrate, and let
$\varepsilon\downarrow0$.  Its initial value is $\varepsilon$;
the stated formula follows.

In particular, the force-to-state map has zero kernel at this based
initial state: $e^\theta=e^u=0$ throughout the interval implies
$h=z=0$ throughout the interval.  It also gives the nonzero target cost
\begin{equation}\label{ifr:force-cost}
 \int_a^t F_\lambda(s)\,ds
 \geq
 \exp\!\left(-\int_a^t C_\lambda(\tau)\,d\tau\right)E_\lambda(t).
\end{equation}
This uses $C_\lambda\geq0$ in \eqref{ifr:reachability-bound}.
It relates the actual signed residual to a reached state, not to the
unsigned termwise majorants rejected earlier.

For the compact affine correction \eqref{ifs:physical-fields},
the target has a further exact lower bound:
\begin{equation}\label{ifr:affine-target-cost}
 E_\lambda(t)\geq
 \frac{\sqrt\pi\,r^2}{8}
 \left(\|L(t)\|_{\mathrm F}^2+\lambda|g(t)|^2\right)^{1/2}.
\end{equation}
To prove it, integrate only over $B_{r/2}$, where
$z=Lx$ and $h=g\cdot x$.  The original coordinate integrals are
\[
 \int_{B_{r/2}}x_i x_j\,dx
    =\delta_{ij}\frac{\pi r^4}{64}.
\]
Thus $\int|Lx|^2=(\pi r^4/64)\operatorname{tr}(L^TL)$
and $\int|g\cdot x|^2=(\pi r^4/64)|g|^2$ on that ball,
which proves the formula with the full Frobenius norm.
Combining \eqref{ifr:force-cost} and
\eqref{ifr:affine-target-cost} gives a quantitative necessary
physical $L^1_tL^2_x$ force budget for the complete compact operation.
Every radius, weight, parent gradient and exponential remains explicit.

This calculation supplies the next test for a proposed annulus
correction: it must satisfy the signed defect requirement
\eqref{ifs:defect-maps} and the complete-state cost
\eqref{ifr:force-cost}.  The estimate neither bounds the successive
parent gradients uniformly nor constructs the next growing return.
Those receiving calculations remain active.  The zero-kernel result
shows why simply replacing the whole increment by a zero-force
evolution from its unchanged initial jets would erase the increment
itself.  A successful construction must retain and estimate the actual
control that reaches its new state.
